\documentclass{lmcs}
\pdfoutput=1

\keywords{choreographic language, PRISM, probabilistic systems modelling}

\usepackage{wrapfig}
\usepackage{sml-highlight}
\usepackage{listings}
\usepackage{xcolor}
\usepackage{comment}
\usepackage{graphicx}
\graphicspath{{./}}
\usepackage{proof}
\usepackage{amssymb}
\usepackage{booktabs}
\usetikzlibrary{automata, positioning, arrows}
\tikzset{every picture/.style={line width=0.75pt}}

\input{macros}

\def\eg{{\em e.g.}}
\def\cf{{\em cf.}}


\begin{document}

\title[A Probabilistic Choreography Language for PRISM]
{A Probabilistic Choreography Language for PRISM}

\author[M.~Carbone]{Marco Carbone\lmcsorcid{0000-0001-9479-2632}}[a]
\author[A.~Veschetti]{Adele Veschetti\lmcsorcid{0000-0002-0403-1889}}[b]

\address{Computer Science Department, IT University of Copenhagen, Rued Langgaards Vej 7, 2300 Copenhagen S, Denmark}	%optional
\email{carbonem@itu.dk, maca@itu.dk}  %optional

\address{Department of Computer Science, TU Darmstadt, Hochschulstraße 10, 64289 Darmstadt, Germany} %optional
\email{adele.veschetti@tu-darmstadt.de} 

\begin{abstract}
  \noindent We present a choreographic framework for modelling and
  analysing concurrent probabilistic systems based on the PRISM
  model-checker. This is achieved through the development of a
  choreography language, which is a specification language that allows
  to describe the desired interactions within a concurrent system from
  a global viewpoint. Using choreographies provides a global view of
  system interactions, which can help in understanding process flow
  and identifying potential modelling issues.  We equip our language
  with a probabilistic semantics and then define a formal encoding
  into the PRISM language and discuss its correctness. Properties of
  programs written in our choreographic language can be model-checked
  by the PRISM model-checker via their translation into the PRISM
  language.  Finally, we implement a compiler for our language and
  demonstrate its practical applicability via examples drawn from use
  cases featured in the PRISM website.
\end{abstract}

\maketitle

\section{Introduction}\label{sec:intro}
\input introduction

\section{A Motivating Example}\label{sec:example}
\input example

\section{Choreography Language}\label{sec:chor}
\input chor

\section{The PRISM Language}\label{sec:prism}
\input prism

\section{Projection}\label{sec:proj}
\input proj

\section{Benchmarking}\label{sec:bench}
\input tests

\section{Related Work and Discussion}\label{sec:relwork}
\input discussion


\bibliographystyle{alphaurl}
\bibliography{biblio}

\end{document}
